% =========================================================
% Proof Stub: Cauchy Sequences are Bounded
% Source: volume-iii/analysis/sequences/notes/sequences/notes-cauchy.tex
% Mode: standard
% =========================================================

\clearpage
\phantomsection
\hypertarget{proof-RA-ABB-P-cauchy-bounded}{}

\subsubsection[Cauchy Sequences are Bounded]{Proof --- Cauchy Sequences are Bounded}

\bigskip

\noindent
\textbf{Source.}
Abbott, \emph{Understanding Analysis}, Lemma~2.6.3.
See notes \texttt{notes-cauchy.tex}.

\vspace{0.75em}

\noindent
\textbf{Goal.}
Every Cauchy sequence is bounded.

\vspace{0.75em}

\noindent
\textbf{Logical form.}
\[
  (x_n) \text{ Cauchy}
  \;\Rightarrow\;
  \exists M > 0,\; \forall n,\; |x_n| \le M.
\]

\vspace{0.75em}

\noindent
\textbf{Proof strategy.}
Use $\varepsilon = 1$: get $N$ so that $|x_m - x_n| < 1$ for $m, n \ge N$.
In particular, $|x_n - x_N| < 1$ for $n \ge N$, so $|x_n| < |x_N| + 1$.
Handle the initial segment $\{x_1, \ldots, x_{N-1}\}$ by taking the max.

\vspace{0.75em}

\noindent
\textbf{Proof.}
\begin{proof}
\Claim Every Cauchy sequence is bounded.

\Given $(x_n)$ a Cauchy sequence.

\Goal Find $M > 0$ with $|x_n| \le M$ for all $n$.

\Strategy Mirroring the proof that convergent implies bounded: use $\varepsilon = 1$
to control the tail; bound the head by finite max.

\medskip

% --- use ε = 1 in Cauchy: get N with |x_m - x_n| < 1 for m,n ≥ N ---

% --- fix n = N: |x_n| < |x_N| + 1 for n ≥ N ---

% --- set M = max(|x_1|, ..., |x_{N-1}|, |x_N| + 1) ---

% --- conclude |x_n| ≤ M for all n ---

\end{proof}

\begin{remark*}[Proof shape]
Same finite-segment trick as ``convergent implies bounded'': control the
tail, take a max over the head.
\end{remark*}

\begin{remark*}[Dependencies]
\begin{itemize}
  \item Definition of Cauchy sequence.
  \item Reverse triangle inequality.
  \item Every finite set of real numbers has a maximum.
\end{itemize}
\end{remark*}
